%% Page 2: Introduction, Figure 2 and the start of Section 2.
\crumbs{§1 Introduction \,›\, Figure~2: GSM8K solve rates \,›\, §2 Chain-of-thought prompting}
\insection{intro,method}

%% ---- marks on the original page
\badge{1}{32.5}{37.4}
\badge{2}{123.5}{33.5}
\bracket{3}{76.0}{137.3}
\bracket{4}{139.5}{177.7}
\bracket{5}{180.0}{214.2}
\bracket{6}{222.2}{254.9}
\highlight[hlpink]{38.1}{58.7}{118.5}{62.2}
\ring[ok]{165.7}{63.9}{169.2}{67.4}
\ring{158.0}{75.4}{161.5}{78.9}
\highlight{90.5}{166.3}{146.4}{170.0}
\penunderline{73.5}{147.0}{146.9}
\highlight[hlgreen]{67.3}{199.2}{122.4}{202.7}

\block[y=34]{%
\num{1}\ \textbf{Why this paper}: bigger models brought better accuracy and sample efficiency, but
\hl[hlpink]{scaling alone \orig{has not proved sufficient}} for multi-step reasoning: arithmetic, commonsense,
symbolic (Rae et al., 2021).}

\block{%
\num{2}\ \textbf{Figure 2}. $y$ = \textbf{solve rate}: \% of GSM8K problems answered correctly.
\begin{smalltable}{l r}
System & solve rate (\%) \\ \midrule
Finetuned GPT-3 175B & 33 \\
Prior best & 55 \\
PaLM 540B, standard prompting & \red{\textbf{18}} \\
PaLM 540B, chain-of-thought prompting & \green{\textbf{57}} \\
\end{smalltable}
\begin{points}
  \item The first two are \textbf{finetuned} results from Cobbe et al. (2021); the last two are the \textbf{same} PaLM
        with prompting only. From 18 to 57 with no training at all.
\end{points}}

\block[y=80]{%
\num{3}\ \textbf{Two ideas combined}:
\begin{steps}
  \item \textbf{Rationales}: models that write natural-language steps before the answer do better at arithmetic,
        but so far they had to be \textbf{trained or finetuned} on such steps.
  \item \textbf{Few-shot prompting} (Brown et al., 2020): no training, just a few input--output
        \orig{exemplars} in the prompt.
\end{steps}}

\block[y=139.5]{%
\num{4}\ \textbf{Why each idea alone falls short}:
\begin{arrows}
  \item rationale training needs a large set of \textbf{high-quality rationales} (underlined): costly to write;
  \item plain few-shot prompting works poorly on reasoning and barely improves with scale.
\end{arrows}
\begin{keyidea}
Prompt with triples \hl{$\langle$input, chain of thought, output$\rangle$}: the benefit of rationales, without training.
\end{keyidea}}

\block[y=181]{%
\num{5}\ Evaluated on \textbf{arithmetic, commonsense and symbolic} reasoning. Why prompting-only matters:
\hl[hlgreen]{no large training dataset}, and \textbf{one} model checkpoint does many tasks.}

\block[y=222]{%
\num{6}\ \textbf{§2 starts from how people solve problems}: ``she has 10 \ldots{} then \ldots{} 7 \ldots{} so the answer
is 7'': split the problem, solve each step, then answer. The goal: give a model the same habit.
\begin{question}
Isn't a chain of thought just an explanation? Not quite: an explanation usually comes \textbf{after} the answer; a
chain of thought comes \textbf{before} it and leads to it (page 3).
\end{question}}
